\documentclass[a4paper,11pt,oneside]{memoir}

% Castellano
\usepackage[spanish,es-tabla]{babel}
\selectlanguage{spanish}
\usepackage[utf8]{inputenc}
\usepackage{placeins}

\RequirePackage{booktabs}
\RequirePackage[table]{xcolor}
\RequirePackage{xtab}
\RequirePackage{multirow}

% Links
\usepackage[colorlinks]{hyperref}
\hypersetup{
	allcolors = {red}
}

% Ecuaciones
\usepackage{amsmath}

% Rutas de fichero / paquete
\newcommand{\ruta}[1]{{\sffamily #1}}

% Párrafos
\nonzeroparskip

% citas
\usepackage{dirtytalk}

% Float barrier
\usepackage{placeins}

% float
\usepackage{float}

% inglés
\newcommand{\english}[1]{\emph{#1}}

% Imagenes
\usepackage{graphicx}
\newcommand{\imagen}[2]{
	\begin{figure}[!h]
		\centering
		\includegraphics[width=0.9\textwidth]{#1}
		\caption{#2}\label{fig:#1}
	\end{figure}
	\FloatBarrier
}

\newcommand{\imagenflotante}[2]{
	\begin{figure}%[!h]
		\centering
		\includegraphics[width=0.9\textwidth]{#1}
		\caption{#2}\label{fig:#1}
	\end{figure}
}



% El comando \figura nos permite insertar figuras comodamente, y utilizando
% siempre el mismo formato. Los parametros son:
% 1 -> Porcentaje del ancho de página que ocupará la figura (de 0 a 1)
% 2 --> Fichero de la imagen
% 3 --> Texto a pie de imagen
% 4 --> Etiqueta (label) para referencias
% 5 --> Opciones que queramos pasarle al \includegraphics
% 6 --> Opciones de posicionamiento a pasarle a \begin{figure}
\newcommand{\figuraConPosicion}[6]{%
  \setlength{\anchoFloat}{#1\textwidth}%
  \addtolength{\anchoFloat}{-4\fboxsep}%
  \setlength{\anchoFigura}{\anchoFloat}%
  \begin{figure}[#6]
    \begin{center}%
      \Ovalbox{%
        \begin{minipage}{\anchoFloat}%
          \begin{center}%
            \includegraphics[width=\anchoFigura,#5]{#2}%
            \caption{#3}%
            \label{#4}%
          \end{center}%
        \end{minipage}
      }%
    \end{center}%
  \end{figure}%
}

%
% Comando para incluir imágenes en formato apaisado (sin marco).
\newcommand{\figuraApaisadaSinMarco}[5]{%
  \begin{figure}%
    \begin{center}%
    \includegraphics[angle=90,height=#1\textheight,#5]{#2}%
    \caption{#3}%
    \label{#4}%
    \end{center}%
  \end{figure}%
}
% Para las tablas
\newcommand{\otoprule}{\midrule [\heavyrulewidth]}
%
% Nuevo comando para tablas pequeñas (menos de una página).
\newcommand{\tablaSmall}[5]{%
 \begin{table}
  \begin{center}
   \rowcolors {2}{gray!35}{}
   \begin{tabular}{#2}
    \toprule
    #4
    \otoprule
    #5
    \bottomrule
   \end{tabular}
   \caption{#1}
   \label{tabla:#3}
  \end{center}
 \end{table}
}

%
% Nuevo comando para tablas pequeñas (menos de una página).
\newcommand{\tablaSmallSinColores}[5]{%
 \begin{table}[H]
  \begin{center}
   \begin{tabular}{#2}
    \toprule
    #4
    \otoprule
    #5
    \bottomrule
   \end{tabular}
   \caption{#1}
   \label{tabla:#3}
  \end{center}
 \end{table}
}

\newcommand{\tablaApaisadaSmall}[5]{%
\begin{landscape}
  \begin{table}
   \begin{center}
    \rowcolors {2}{gray!35}{}
    \begin{tabular}{#2}
     \toprule
     #4
     \otoprule
     #5
     \bottomrule
    \end{tabular}
    \caption{#1}
    \label{tabla:#3}
   \end{center}
  \end{table}
\end{landscape}
}

%
% Nuevo comando para tablas grandes con cabecera y filas alternas coloreadas en gris.
\newcommand{\tabla}[6]{%
  \begin{center}
    \tablefirsthead{
      \toprule
      #5
      \otoprule
    }
    \tablehead{
      \multicolumn{#3}{l}{\small\sl continúa desde la página anterior}\\
      \toprule
      #5
      \otoprule
    }
    \tabletail{
      \hline
      \multicolumn{#3}{r}{\small\sl continúa en la página siguiente}\\
    }
    \tablelasttail{
      \hline
    }
    \bottomcaption{#1}
    \rowcolors {2}{gray!35}{}
    \begin{xtabular}{#2}
      #6
      \bottomrule
    \end{xtabular}
    \label{tabla:#4}
  \end{center}
}

%
% Nuevo comando para tablas grandes con cabecera.
\newcommand{\tablaSinColores}[6]{%
  \begin{center}
    \tablefirsthead{
      \toprule
      #5
      \otoprule
    }
    \tablehead{
      \multicolumn{#3}{l}{\small\sl continúa desde la página anterior}\\
      \toprule
      #5
      \otoprule
    }
    \tabletail{
      \hline
      \multicolumn{#3}{r}{\small\sl continúa en la página siguiente}\\
    }
    \tablelasttail{
      \hline
    }
    \bottomcaption{#1}
    \begin{xtabular}{#2}
      #6
      \bottomrule
    \end{xtabular}
    \label{tabla:#4}
  \end{center}
}

%
% Nuevo comando para tablas grandes sin cabecera.
\newcommand{\tablaSinCabecera}[5]{%
  \begin{center}
    \tablefirsthead{
      \toprule
    }
    \tablehead{
      \multicolumn{#3}{l}{\small\sl continúa desde la página anterior}\\
      \hline
    }
    \tabletail{
      \hline
      \multicolumn{#3}{r}{\small\sl continúa en la página siguiente}\\
    }
    \tablelasttail{
      \hline
    }
    \bottomcaption{#1}
  \begin{xtabular}{#2}
    #5
   \bottomrule
  \end{xtabular}
  \label{tabla:#4}
  \end{center}
}



\definecolor{cgoLight}{HTML}{EEEEEE}
\definecolor{cgoExtralight}{HTML}{FFFFFF}

%
% Nuevo comando para tablas grandes sin cabecera.
\newcommand{\tablaSinCabeceraConBandas}[5]{%
  \begin{center}
    \tablefirsthead{
      \toprule
    }
    \tablehead{
      \multicolumn{#3}{l}{\small\sl continúa desde la página anterior}\\
      \hline
    }
    \tabletail{
      \hline
      \multicolumn{#3}{r}{\small\sl continúa en la página siguiente}\\
    }
    \tablelasttail{
      \hline
    }
    \bottomcaption{#1}
    \rowcolors[]{1}{cgoExtralight}{cgoLight}

  \begin{xtabular}{#2}
    #5
   \bottomrule
  \end{xtabular}
  \label{tabla:#4}
  \end{center}
}


















\graphicspath{ {./img/} }

% Capítulos
\chapterstyle{bianchi}
\newcommand{\capitulo}[2]{
	\setcounter{chapter}{#1}
	\setcounter{section}{0}
	\chapter*{#2}
	\addcontentsline{toc}{chapter}{#2}
	\markboth{#2}{#2}
}

% Apéndices
\renewcommand{\appendixname}{Apéndice}
\renewcommand*\cftappendixname{\appendixname}

\newcommand{\apendice}[1]{
	%\renewcommand{\thechapter}{A}
	\chapter{#1}
}

\renewcommand*\cftappendixname{\appendixname\ }

% Formato de portada
\makeatletter
\usepackage{xcolor}
\newcommand{\tutor}[1]{\def\@tutor{#1}}
\newcommand{\course}[1]{\def\@course{#1}}
\definecolor{cpardoBox}{HTML}{E6E6FF}
\def\maketitle{
  \null
  \thispagestyle{empty}
  % Cabecera ----------------
\noindent\includegraphics[width=\textwidth]{cabecera}\vspace{1cm}%
  \vfill
  % Título proyecto y escudo informática ----------------
  \colorbox{cpardoBox}{%
    \begin{minipage}{.8\textwidth}
      \vspace{.5cm}\Large
      \begin{center}
      \textbf{TFG del Grado en Ingeniería Informática}\vspace{.6cm}\\
      \textbf{\LARGE\@title{}}
      \end{center}
      \vspace{.2cm}
    \end{minipage}

  }%
  \hfill\begin{minipage}{.20\textwidth}
    \includegraphics[width=\textwidth]{escudoInfor}
  \end{minipage}
  \vfill
  % Datos de alumno, curso y tutores ------------------
  \begin{center}%
  {%
    \noindent\LARGE
    Presentado por \@author{}\\ 
    en Universidad de Burgos --- \@date{}\\
    Tutores: \@tutor{}\\
  }%
  \end{center}%
  \null
  \cleardoublepage
  }
\makeatother

\newcommand{\nombre}{Daniel Arnaiz Gutierrez} %%% cambio de comando

% Datos de portada
\title{Extracción y procesamiento de datos de Amazon para su utilización en un estudio de marketing.}
\author{\nombre}
\tutor{Dr. José Francisco Díez Pastor \\ Dr. César Ignacio García Osorio}
\date{\today}

\begin{document}

\maketitle


\newpage\null\thispagestyle{empty}\newpage


%%%%%%%%%%%%%%%%%%%%%%%%%%%%%%%%%%%%%%%%%%%%%%%%%%%%%%%%%%%%%%%%%%%%%%%%%%%%%%%%%%%%%%%%
\thispagestyle{empty}


\noindent\includegraphics[width=\textwidth]{cabecera}\vspace{1cm}

\noindent Dr. José Francisco Díez Pastor y Dr. César Ignacio García Osorio, profesores del departamento de Ingeniería Civil, área de Lenguajes y sistemas informáticos.

\noindent Exponen:

\noindent Que el alumno D. \nombre, con DNI 71296333, ha realizado el Trabajo final de Grado en Ingeniería Informática titulado <<Extracción y procesamiento de datos de Amazon para su utilización en un estudio de marketing>>. 

\noindent Y que dicho trabajo ha sido realizado por el alumno bajo la dirección del que suscribe, en virtud de lo cual se autoriza su presentación y defensa.

\begin{center} %\large
En Burgos, {\large \today}
\end{center}

\vfill\vfill\vfill

% Author and supervisor
\begin{minipage}{0.45\textwidth}
\begin{flushleft} %\large
Vº. Bº. del tutor:\\[2cm]
D. José Francisco Díez Pastor
\end{flushleft}
\end{minipage}
\hfill
\begin{minipage}{0.45\textwidth}
\begin{flushleft} %\large
Vº. Bº. del tutor:\\[2cm]
D. César Ignacio García Osorio
\end{flushleft}
\end{minipage}
\hfill

\vfill

% para casos con solo un tutor comentar lo anterior
% y descomentar lo siguiente
%Vº. Bº. del Tutor:\\[2cm]
%D. nombre tutor


\newpage\null\thispagestyle{empty}\newpage




\frontmatter

% Abstract en castellano
\renewcommand*\abstractname{Resumen}
\begin{abstract}
El principal objetivo de este proyecto es la recolección de datos de determinados productos del sitio web <<Amazon.com>> para después procesar y almacenar los resultados de forma relacional. Además, parte de este proceso consiste en, dada una imagen del producto, etiquetarla de forma automática según las características de esta.

Para realizar la parte de la recolección de datos se utilizará una técnica de minería de datos llamada \english{web scraping} y para el procesado y etiquetado de las imágenes se harán uso de redes neuronales, otra rama de la minería de datos.

Los resultados, tanto del \english{scraper} como del clasificador de imágenes, serán almacenados en una base de datos para su posterior uso en estudios de marketing.
\end{abstract}

\renewcommand*\abstractname{Descriptores}
\begin{abstract}
Minería de datos, \english{Web scraping}, Clasificador, Base de datos, Redes neuronales, Inteligencia artificial
\end{abstract}

\clearpage

% Abstract en inglés
\renewcommand*\abstractname{Abstract}
\begin{abstract}
The main objective of this project is the data collection of certain products in the website <<Amazon.com>> to later on process and store the result in a relational way. In addition to this, part of this process consists of, given a product image, label it automatically depending of its characteristics.

To perform the part in which the data is collected, a data mining technique called web scraping will be used. Also, for the processing and labeling of the images, neuronal networks will be used, which is another branch of data mining.

The results, both of the scraper and the image classifier, will be stored in a database for later use in marketing studies.

\end{abstract}

\renewcommand*\abstractname{Keywords}
\begin{abstract}
Data mining, Web scraping, Classifier, Database, Neuronal network, Artificial intelligence
\end{abstract}

\clearpage

% Indices
\tableofcontents

\clearpage

\listoffigures

\clearpage

%\listoftables
%\clearpage

\mainmatter
\capitulo{1}{Introducción}

Uno de los principales problemas que se encuentran a la hora de hacer un estudio de mercado es el conjunto de datos sobre el que trabajar. Existen muchas maneras de recolectar esos datos como podrían ser encuestas, sondeos, entrevistas o incluso observaciones directas. Es aquí donde entra el proceso de recolección de datos usado en este proyecto: \english{Web scraping}. Este concepto se puede definir como una técnica de obtención de datos de sitios web de forma automatizada, forma parte de un concepto que lo engloba llamado \english{Data mining} o Minería de datos.

Este proceso conlleva una serie de ventajas sobre la recolección de datos de forma manual:

\begin{itemize}
	\item Una de las principales razones por la que utilizar un \english{web scraper} es la automatización, dado que siempre vas a querer extraer los mismos campos o características de un sitio web, este proceso se convierte rápidamente en repetitivo y aburrido. Un \english{web scraper} es capaz de realizar estas acciones de forma relativamente autónoma y mucho más rápido de lo que sería hacerlo de forma manual.
	\item Otra ventaja importante de utilizar un \english{web scraper} sobre la extracción de datos de forma manual sería el almacenamiento de la información extraída. Aunque se podrían llegar a almacenar de la misma forma los datos extraídos, con un \english{web scraper} tienes la posibilidad de personalizar el modo en que se almacenen los datos, de tal forma que sea una propia extensión del proceso de extracción, además de poder guardarlo en varios formatos a la vez sin apenas alargar la duración del proceso.
\end{itemize}

Una vez que se dispone de la información necesaria para iniciar el estudio, es probable que se necesite de un procesamiento de esta información para clasificar o identificar los datos deseados. Existen infinidad de procesamientos de datos posibles, desde procesamiento de textos como \english{Text sentiment analysis}\footnote{Proceso automático de identificación y extracción de información subjetiva de un texto.} o categorización, a procesamiento de imágenes como se ha hecho en este proyecto. En este caso, se ha hecho una clasificación automática de imágenes para más tarde extraer información de éstas dependiendo de la categoría en la que han sido clasificadas.
\capitulo{2}{Objetivos del proyecto}

Este apartado explica de forma precisa y concisa cuales son los objetivos que se persiguen con la realización del proyecto. Se puede distinguir entre los objetivos marcados por los requisitos del software a construir y los objetivos de carácter técnico que plantea a la hora de llevar a la práctica el proyecto.

\section{Objetivos de carácter general}

El principal objetivo de este proyecto es la creación de varias herramientas capaces de extraer y procesar información sobre un familia de productos concreta de Amazon, en este caso camisetas, para su posterior uso en estudios de mercado.

A su vez, este proceso está divido en diferentes fases que podemos identificar como objetivos secundarios:

\begin{itemize}
	\item Ser capaz de extraer información de los distintos productos. Esta información incluye campos como el nombre de la marca, el rango de precios, número de valoraciones o imágenes del producto. Este objetivo se realiza con técnicas de \english{web scraping}.
	\item Como complemento al anterior objetivo, otro objetivo es la implementación de un clasificador automático de imágenes. Este clasificador etiquetará cada imagen en función de si aparece un modelo o no, y si la cara del modelo es visible o no. Esto se consigue con técnicas de inteligencia artificial y redes neuronales, y más concretamente con el uso de \english{deep learning}.
	\item Otro de los principales objetivos del proyecto es el almacenamiento de la información descargada y de las imágenes clasificadas. Para cumplir con este objetivo se hará uso de un gestor de bases de datos y se almacenará la información en múltiples tablas con sus respectivos campos.
	\item Por último, un objetivo impuesto para la realización de este proyecto ha sido la visualización de la información descargada y procesada en un documento Excel.
\end{itemize}

\section{Objetivos de carácter técnico}

A continuación se listan los objetivos técnicos utilizados para la correcta implementación de los objetivos de carácter mas generales expuestos anteriormente.

\begin{itemize}
	\item Utilizar Python como lenguaje de programación para la totalidad del proyecto.
	\item Hacer uso de la herramienta Scrapy para implementar todo lo relacionado con la extracción de datos y \english{Web scraping}.
	\item Un objetivo técnico importante y necesario para la creación del clasificador automático es la creación de un \english{dataset} o conjunto de datos personalizado de forma manual para posteriormente entrenar el clasificador. Para ello se ha usado la herramienta Dataturks\footnote{\url{https://dataturks.com/}}.
	\item Usar JSON como forma de almacenamiento temporal de los productos descargados debido la simplicidad de su estructura.
	\item Utilizar SQLite\footnote{\url{https://www.sqlite.org}} como gestor de base de datos para el almacenamiento final de los productos descargados.
	\item Implementar el clasificador de imágenes con la ayuda de la librería Keras\footnote{\url{https://keras.io/}} y de Google Colab\footnote{\url{https://colab.research.google.com/}}.
	\item Utilizar el servicio de GitHub\footnote{\url{https://github.com/}} para hacer uso de un control de versiones del proyecto.
	\item Seguir una metodología ágil para el desarrollo y el enfoque del proyecto. Para esto se ha usado la herramienta ZenHub\footnote{\url{https://www.zenhub.com/}} como una extensión de GitHub.
\end{itemize}
\capitulo{3}{Conceptos teóricos}

\section{Inteligencia artificial}

Un concepto tan amplio como lo es la inteligencia artificial se puede definir de múltiples formas, a continuación se citan algunas:

\begin{enumerate}
    \item \say{La habilidad, como sistema, de interpretar correctamente datos externos, aprender de estos datos, y usar lo aprendido para lograr objetivos y tareas concretas a través de una adaptación flexible} \cite{kaplan_haenlein_2019}.
    \item \say{... la habilidad de un sistema de actuar adecuadamente en un ecosistema incierto en el que la acción apropiada es la que aumenta la probabilidad de alcanzar el objetivo, siendo el objetivo completar una serie de logros menores que ayuden a alcanzar el objetivo final} \cite{albus_1991}.
    \item \say{Estudio de del diseño de agentes inteligentes} \cite{poole_mackworth_goebel_1998}.
\end{enumerate}

\subsection{Historia de la inteligencia artificial}
Durante la segunda guerra mundial, el aclamado científico Alan Turing descifró el <<Código Enigma>> usado por las fuerzas alemanas para mandar mensajes de forma segura. Alan Turing y Enigma conforman los pilares de la inteligencia artificial y el \english{machine learning}.

En 1956, el científico John McCarthhy\footnote{\url{https://es.wikipedia.org/wiki/John\_McCarthy}}, usó el término inteligencia artificial por primera vez en la Conferencia Dartmouth\footnote{\url{https://es.wikipedia.org/wiki/Conferencia\_de\_Dartmouth}}, dando lugar a la creación de varios centros de investigación en Estados Unidos con el propósito del estudio del potencial de la inteligencia artificial. Los investigadores Allen Newell\footnote{\url{https://es.wikipedia.org/wiki/Allen\_Newell}} y Herbert Simon\footnote{\url{https://es.wikipedia.org/wiki/Herbert\_Alexander\_Simon}}  fueron clave para impulsar inteligencia artificial como un nuevo campo en la Informática que podría cambiar el mundo.

A partir de los años 60, el principal enfoque de los investigadores se centró en desarrollar algoritmos para resolver problemas y teoremas matemáticos. Este proceso dio lugar a los primeros conceptos sobre \english{machine vision learning}\footnote{\url{https://es.wikipedia.org/wiki/Vision\_artificial}} y inteligencia artificial aplicada a la robótica.

Tras un largo periodo sin avances notables debido a una escasez de financiaciones para el desarrollo de la inteligencia artificial, a finales de los años 90, el desarrollo resurge gracias a grandes corporaciones interesadas en usar este campo para lograr unas ventajas y beneficios nunca vistos, Es así como en 1997 la máquina Deep Blue\footnote{\url{https://es.wikipedia.org/wiki/Deep_Blue_(computadora)}} de IBM se convirtió en la primera computadora capaz de derrotar a un campeón mundial de ajedrez como lo fue Garry Kasparov.

En los ultimos 15 años, las principales empresas mundiales han utilizado la inteligencia artificial como una herramienta más a la hora de crecer y convertirse en el sector más importante del mundo, gracias a campos como la visión artificial, minería de datos o el procesamiento de lenguajes naturales \cite{aihistory}.

\subsection{Machine learning (Aprendizaje automático)}
El \english{Machine learning} se considera un campo de las ciencias de computación y a su vez una rama de la inteligencia artificial. Tiene como objetivo principal el desarrollo de técnicas y la búsqueda de algoritmos y heurísticas que permitan que las computadoras aprendan por si mismas a través de la experiencia.

Las aplicaciones que conforman este campo son muy variadas, desde su uso en diagnósticos médicos, análisis de patrones, reconocimiento de lenguajes naturales hablados y escritos, a juegos o robótica \cite{wiki:apredizaje_automatico}.

A su vez los diferentes algoritmos de aprendizaje automático se pueden agrupar por tipos según la salida de los mismos.
Encontramos así:

\begin{description}
    \item[Aprendizaje supervisado] Produce una función que establece una correspondencia entre las entradas y salidas deseadas. Un posible uso de estos algoritmos podría ser un problema de clasificación en el que hay que etiquetar una serie de valores usando una entre varias categorías. En este tipo de aprendizaje, la base del conocimiento estaría formada por ejemplos de etiquetados anteriores.
    \item[Aprendizaje no supervisado] En este caso el proceso se lleva a cabo por un conjunto de datos únicamente formado por entradas al sistema que no contienen información sobre las categorías de cada ejemplo. El sistema debe reconocer los patrones de forma autónoma para encontrar grupos de elementos comunes y asignar las nuevas entradas a dichos grupos.
    \item[Aprendizaje semi-supervisado] Estos algoritmos combinan los dos anteriores para clasificar teniendo en cuenta los datos clasificados y los no clasificados.
    \item[Aprendizaje por refuerzo] En este tipo de algoritmos, el sistema aprende observando lo que le rodea, es decir, su información de entrada se basa en la retroalimentación que obtiene como respuesta a sus acciones. Se podría decir que aprende a base de prueba-error.
    \item[Transducción] Un algoritmo similar al aprendizaje supervisado tratado anteriormente. La principal diferencia sería que no construye de forma explícita una función. Un posible ejemplo de este tipo de algoritmos sería el $k$ de vecinos mas cercanos\footnote{\url{https://es.wikipedia.org/wiki/K\_vecinos\_mas\_proximos}}.
    \item[Aprendizaje multi-tarea] Los métodos incluidos en esta categoría se basan en la utilización de el conocimiento aprendido previamente para enfrentarse a problemas similares a los ya vistos.
\end{description}

\section{Redes neuronales}
Una Red Neuronal Artificial (RNA) se basa en un modelo inspirado en el comportamiento de las neuronas y como se organizan para crear una estructura en el cerebro.

El cerebro se puede llegar a ver como un sistema inteligente, aunque de manera distinta a como trabajan los ordenadores hoy en día. Si bien las computadoras actuales son muy rápidas y eficientes en el procesamiento de información, existen otro tipo de tareas y procesos mas complejos que requieren más tiempo y más esfuerzo computacional que a un cerebro humano. Este tipo de tareas están relacionadas con el reconocimiento y clasificación de patrones \cite{redes_neuronales}.

\subsection{Neurona artificial}
El primer modelo de neurona artificial, con el objetivo de llevar a cabo tareas muy simples, fue creado por el neuroanatomista Warren McCulloch y el matemático Walter Pitts y consta de las siguientes partes:
\begin{itemize}
    \item Conjunto de entradas $x_1$, $x_2$, ... $x_n$
    \item Pesos sinápticos $w_1$, $w_2$, ... $w_n$, correspondientes a cada entrada
    \item Una función de agregación, $\sum$
    \item Una función de activación, $f$
    \item Una salida, $Y$
\end{itemize}
\imagen{neurona_artificial}{Diferenciación entre una neurona biológica y una neurona artificial}
%\caption[Diferenciación entre una neurona biológica y una neurona artificial]{Extraído de \url{http://www.cs.us.es/\~fsancho/?e=72}}

\subsection{Perceptrón}
 Un perceptrón es el modelo más simple de una neurona artificial, está basado en una única neurona y posee una única salida. El perceptrón puede combinarse con otros tipos de perceptrones o neuronas artificiales para formar una red neuronal mas compleja.
 
 A su vez, un perceptrón puede organizar sus neuronas en capas: capa de entrada, una capa de salida, y múltiples capas ocultas formadas por neuronas. En el caso de existir capas ocultas nos referimos a la red neuronal como perceptrón multicapa\footnote{\url{https://es.wikipedia.org/wiki/Perceptron\_multicapa}}

\imagen{perceptron_multicapa}{Diagrama de un  perceptrón multicapa}
%\captionDiagrama de un  perceptrón multicapa]{Extra ido de \url{https://es.wikipedia.org/wiki/Perceptron\_multicapa}}

\section{Minería de datos}

La minería de datos se considera un campo de la estadística y de las ciencias de la computación. El proceso de la minería de datos se basa en el descubrimiento de patrones en grandes conjuntos de datos utilizando métodos de inteligencia artificial, aprendizaje automático, estadística y bases de datos. El principal objetivo es el de extraer información de un conjunto de datos y procesarlo para transformarlo de forma estructurada y comprensible para su uso posterior \cite{wiki:mineria_de_datos}.

El proceso típico de minería de datos está constituido por los siguientes pasos generales:

\begin{description}
    \item[Selección del \english{dataset} o conjunto de datos] Selección de los campos que se quieren predecir y de los necesarios para el proceso.
    \item[Análisis de las propiedades de los datos] Comprobación de histogramas, diagramas de dispersión y búsqueda de valores atípicos o nulos.
    \item[Transformación del conjunto de datos de entrada] Preparación de los datos obtenidos para aplicar la técnica de minería de datos que mejor se adapte al problema a resolver.
    \item[Selección y aplicación de la técnica de minería de datos] Construcción del modelo a aplicar: predictivo, o de clasificación.
    \item[Extracción de datos] Proceso por el que se recopilan los datos a ser tratados, ya sea registrando las medidas de sensores, consultando bases de datos, o como en este proyecto, realizando tareas de \english{web scraping}.
    \item[Interpretación y evaluación] Una vez obtenido el modelo, se procede a su validación y evaluación.
\end{description}

\subsection{Técnicas de minería de datos}
\begin{description}
    \item[Redes neuronales] Tal y como se han descrito anteriormente, esta técnica de procesamiento automático y los diferentes tipos de red neuronal (perceptrón y perceptrón multicapa) son una técnica utilizada en la minería de datos.
    \item[Regresión lineal] Técnica utilizada para identificar relaciones lineales entre dos variables.
    \imagen{regresion_lineal}{Ejemplo de diagrama de regresión lineal}
    \item[Árboles de decisión] Modelo de predicción utilizado en el campo de la inteligencia artificial y el análisis predictivo\footnote{\url{https://es.wikipedia.org/wiki/Analisis\_predictivo}} donde a partir de un conjunto de datos se construyen diagramas para representar y categorizar condiciones de forma sucesiva hasta resolver el problema dado.
    \imagen{arbol_decision}{Ejemplo de árbol de decisión}
    \item[Modelos estadísticos]Expresión o ecuación utilizada con el fin de indicar diferentes factores que modifican la variable de una respuesta.
    \item[\english{Clustering}] Proceso de agrupación de una series de valores según un criterio de distancia o similitud.
    \imagen{clustering}{Ejemplo de diagrama mostrando \english{clustering}}
    \item[Reglas de asociación] Procedimiento utilizado para reconocer hechos en común dentro de un conjunto de datos.
\end{description}

\section{Web Scraping}
El web \english{scraping} es un a técnica utilizada para extraer información desde páginas web a través de uno o varias herramientas de software, estos programas simulan la navegación por el sitio web a la hora de extraer los datos deseados. La técnica de web \english{scraping} es muy utilizada a día de hoy por servicios como motores de búsqueda ya que es necesario para indexar la web, sin embargo, los principales usos de esta técnica son la transformación de datos sin estructurar que podemos encontrar en cualquier web en formato HTML\footnote{\url{https://es.wikipedia.org/wiki/HTML}} a estructuras de datos que pueden ser almacenados en bases de datos, hojas de cálculo, o en cualquier forma de almacenamiento estructurada \cite{wiki:web_scraping}.

En cuanto a las técnicas utilizadas para realizar \english{web scraping} sobre un sitio web encontramos varios tipos en función del nivel de automatización del proceso:

\begin{description}
    \item[Expresiones regulares] Es posible extraer información de una página web utilizando expresiones regulares.
    \item[Peticiones HTTP] Se basa en extraer información por medio de peticiones HTTP al servidor remoto utilizando sockets.
    \item[Algoritmos de minería de datos] Aplicación de algunas de las técnicas de minería de datos tratadas en el apartado anterior.
    \item[Parser HTML] Existen lenguajes capaces de procesar documentos HTML y extraer de forma estructurada todo el contenido de una web.
    \item[Aplicaciones de web scraping] Programas completos de web scraping que combinan algunas de las técnicas anteriores para ofrecer un proceso personalizado para cada caso, facilitando y simplificando la obtención del código de la aplicación de obtención de datos a partir de páginas web.
\end{description}
\capitulo{4}{Técnicas y herramientas}

Se han utilizado diferentes herramientas para la realización de este proyecto, a continuación se listan con una pequeña descripción de cada una:


\section{Lenguajes y bibliotecas de programación}


\subsection{Python}
Python\footnote{\url{https://www.python.org/}} es un lenguaje de programación nacido en los años 90 y ampliamente utilizado en el presente. Se trata de un lenguaje interpretado, dinámicamente tipado y multiplataforma. Soporta programación orientada a objetos y sigue una filosofía de legibilidad y transparencia \cite{wiki:python}.

Ha sido el lenguaje utilizado en la mayor parte de este proyecto, desde la parte de \english{web scraping} a la clasificación de imágenes.


\subsection{Scrapy}
Scrapy\footnote{\url{https://scrapy.org/}} es una herramienta gratuita y de código abierto escrita en Python por Scrapinghub Ltd. y originalmente diseñada para \english{web scraping} aunque también utilizado para extraer datos mediante APIs \cite{wiki:scrapy}.

Scrapy está compuesto por una serie de características que lo hacen la herramienta más adecuada para este proyecto. Entre ellas encontramos:

\begin{itemize}
	\item Soporte para extraer y seleccionar datos de fuentes HTML/XML usando selectores CSS\footnote{\url{https://es.wikipedia.org/wiki/Hoja_de_estilos_en_cascada}} y expresiones XPath\footnote{\url{https://es.wikipedia.org/wiki/XPath}}.
	\item Soporte para exportar los datos extraídos en múltiples formatos tales como JSON\footnote{\url{https://es.wikipedia.org/wiki/JSON}}, CSV\footnote{\url{https://es.wikipedia.org/wiki/Valores_separados_por_comas}} y XML\footnote{\url{https://es.wikipedia.org/wiki/Extensible_Markup_Language}}.
	\item Fuerte extensibilidad, ya que permite conectar funcionalidades externas tales como bases de datos.
\end{itemize}


\subsection{Keras}
Keras\footnote{\url{https://keras.io/}} se trata de una librería de redes neuronales de código abierto escrita en Python por François Chollet y capaz de ser ejecutada sobre TensorFlow. Las principales características que lo componen son su modularidad, extensibilidad y el ser amigable con el usuario. Además del soporte para las redes neuronales estándar, Keras ofrece soporte para las Redes Neuronales Convolucionales y Redes Neuronales Recurrentes \cite{wiki:keras}.


\subsection{OpenCV}
OpenCV\footnote{\url{https://opencv.org/}} es una biblioteca multiplataforma originalmente desarrollada por Intel y orientada a la visión artificial. Nace en 1999 y actualmente es utilizada en infinidad de aplicaciones. Contiene más de 500 funciones abarcando una gran gama de áreas en el proceso de visión, como reconocimiento de objetos (reconocimiento facial), calibración de cámaras y visión robótica \cite{wiki:opencv}.


\subsection{SQLite}
SQLite\footnote{\url{https://www.sqlite.org/index.html}} es un sistema de gestión de bases de datos relacional multiplataforma que forma parte de una biblioteca escrita en C por Richard Hipp. Su principal característica es que la base de datos se almacena en un único archivo.


\section{Herramientas de desarrollo}


\subsection{Visual Studio Code}
Visual Studio Code es un editor de código desarrollado por Microsoft anunciado en 2015 y lanzado al público en 2016. Cuenta con soporte para multitud de lenguajes de programación y entre sus características mas destacadas encontramos soporte para depuración, control integrado de Git\footnote{\url{https://git-scm.com/}}, resaltado de sintaxis, refactorización de código y  finalización inteligente de código \cite{wiki:vscode}.

Este editor ha sido utilizado para escribir la mayoría de código que encontramos en el proyecto.


\subsection{Google Colab (Colaboratory)}
Google Colab\footnote{\url{https://colab.research.google.com}} es un entorno de Jupyter Notebook\footnote{\url{https://jupyter.org/}} ya configurado y completamente ejecutado en remoto. Permite editar y ejecutar código, guardar y compartir proyectos y tener acceso a potentes recursos desde un navegador web.

Esta herramienta ha sido usada a la hora de desarrollar y más tarde entrenar el clasificador de imágenes utilizado en el proyecto.


\section{Herramientas de gestión de código}


\subsection{Git}
 Git\footnote{\url{https://git-scm.com/}} es un sistema de control de versiones gratuito y \english{open source} diseñado para gestionar proyectos de forma sencilla y eficiente.

\subsection{GitHub}
GitHub\footnote{\url{https://github.com/}} es una plataforma de desarrollo colaborativo utilizado para alojar proyectos y utilizar el sistema de control de versiones Git \cite{wiki:github}.

\subsection{ZenHub}
ZenHub\footnote{\url{https://www.zenhub.com/}} es una herramienta utilizada para gestión ágil de proyectos. Se integra de forma nativa en GitHub añadiendo potentes características de rastreo y planificación sin necesidad de configuración.

\section{Herramientas de documentación}

\subsection{\LaTeX{}}
\LaTeX{}\footnote{\url{https://es.wikipedia.org/wiki/LaTeX}} es un sistema de composición de textos orientado a la creación de documentos con una alta calidad tipográfica como podrían ser artículos científicos u otros textos con, por ejemplo, expresiones matemáticas.

Se trata de un software libre basado en un conjunto de macros de \TeX, que a su vez se trata de un lenguaje escrito por Leslie Lamport en 1984 \cite{wiki:latex}.

\subsection{Overleaf}
Overleaf\footnote{\url{https://es.overleaf.com/}} es un editor en línea de documentos \LaTeX{} muy intuitivo y fácil de usar. No requiere de instalación y ofrece diversas características como colaboración en tiempo real, control de versiones o el uso de plantillas \LaTeX{}.

\section{Otras técnicas y herramientas}

\subsection{JSON}
JSON (JavaScript Object Notation) es un formato de texto utilizado para el intercambio de datos que desde el año 2019 es considerado como un formato independiente de lenguaje. La principal característica de este lenguaje y la razón por la que se ha utilizado en este proyecto es la facilidad de extraer los datos que lo conforman  con ayuda de \english{parsers} o analizadores sintácticos.

\subsection{Dataturks}
Dataturks\footnote{\url{https://dataturks.com/}} es una herramienta web cuya finalidad es anotar imágenes de forma manual con el propósito de crear conjuntos de entrenamiento para modelos de redes neuronales como Keras o TensorFlow\footnote{\url{https://www.tensorflow.org/}}.

En este proyecto, esta herramienta ha sido indispensable para crear el modelo que clasifica automáticamente las imágenes de cada producto.

















\capitulo{5}{Aspectos relevantes del desarrollo del proyecto}

Este apartado recoge los aspectos más interesantes del desarrollo del proyecto, el ciclo de vida utilizado y los detalles de mayor relevancia de las fases de análisis, diseño e implementación.

\section{Metodologías aplicadas}

Para la realización de este proyecto se ha decidido seguir una metodología ágil siempre que ha sido posible. Por ello, se han seguido ciertas líneas de trabajo. A continuación se listan:

\begin{itemize}
    \item Se ha dividido el desarrollo del proyecto en múltiples \english{sprints}.
    \item Se han ido proponiendo nuevos requisitos y objetivos a medida que se iban cumpliendo los anteriores.
    \item Se han realizado reuniones entre todos los integrantes al acabar cada \english{sprint}.
    \item Al acabar cada \english{sprint} se ha realizado una demostración siempre que los cambios eran sustanciales.
    \item Se han organizado las diferentes tareas y objetivos a completar en un tablero Kanban\footnote{\url{https://kanbantool.com/es/tablero-kanban}}.
    \item Se ha asociado una duración estimada a cada tarea antes de realizarla.
\end{itemize}

\section{Propuesta del proyecto}
La idea de este proyecto nace con el objetivo de crear una herramienta que ayude a obtener y descargar información sobre una familia de productos concreta de la web de compras por internet \english{Amazon.com} con el objetivo de hacer múltiples estudios sobre los diferentes conjuntos de datos obtenidos.

Desde un primer momento se ha querido obtener la mayor cantidad de campos de información, es por ello que una vez el proceso ha sido completado, obtenemos la siguiente información de cada producto:

\begin{description}
    \item[ASIN:] Código identificativo de cada producto.
    \item[Sexo:] Sexo al que va dirigido cada artículo.
    \item[Rango de precios:] El precio de cada producto puede variar en función del color o la talla. Este campo recoge el precio máximo y mínimo de cada producto.
    \item[Marca:] El nombre de la marca de cada artículo.
    \item[Descripción:] Texto que el vendedor ha utilizado para describir y dar detalles de cada producto.
    \item[Puntuación:] Media de las puntuaciones que los clientes han votado a cada producto. Va del 0 al 5.
    \item[Número de valoraciones:] Número de clientes que han proporcionado una valoración del producto.
    \item[Comentarios sobre el producto:] Comentarios completos sobre un producto en el que los compradores dejan su opinión. El número de comentarios a extraer debe ser definido al comenzar el proceso.
    \item[URL:] Dirección web de cada producto
    \item[Imágenes:] Conjunto de imágenes asociadas a cada artículo.
    \item[Contenido de la imagen principal:] Este campo describe el contenido de la imagen principal de cada producto en función de si aparece un modelo o no, y si la cara del modelo es visible o no.
\end{description}

La mayoría de estos campos se pueden extraer directamente del código \english{HTML} de la web de cada producto, pero para ciertos casos es necesario procesar antes la información ya extraída y usarla para crear los campos que faltan. Como es el caso del campo que recoge el contenido de la imagen principal.

A continuación se irán presentando y detallando las diferentes fases de las que se compone el proceso de extracción y procesamiento de datos.

\section{Extracción de los datos}

La herramienta utilizada en este proyecto para realizar el \english{web scraping} ha sido Scrapy. Se han barajado otras opciones como BeautifulSoup\footnote{\url{https://www.crummy.com/software/BeautifulSoup/bs4/doc/}} o Selenium\footnote{\url{https://www.seleniumhq.org/}}.

Las principales diferencias y ventajas que lo separan del resto son:
\begin{itemize}
    \item Es más rápido ya que trabaja de forma asíncrona.
    \item Tiene mejor soporte para parseado de HTML.
    \item Soporta caracteres unicode.
    \item Soporta redirecciones HTTP.
    \item Exportación directa de los resultados a JSON, CSV o XML.
    \item Se puede usar junto a otras librerías como urllib\footnote{\url{https://docs.python.org/3/library/urllib.html}} o BeautifulSoup.
\end{itemize}

\subsection{Arquitectura de Scrapy}

En cuanto a la arquitectura de Scrapy, hay que destacar los tres principales bloques que lo conforman:

\begin{description}
    \item[\english{Items}:] Los diferentes elementos o campos a extraer.
    \item[\english{Spiders}:] Clases escritas por el usuario que contienen los diferentes procedimientos para parsear y extraer los \english{items} de uno o varios enlaces. En este caso hemos hecho uso de 3 \english{spiders}: Una para obtener los enlaces a cada artículo, otra para extraer los comentarios, y la última para el resto de campos.
    \item[\english{Pipelines}:] Apartado responsable de procesar los \english{items} una vez han sido extraidos por las \english{spiders}. En este proyecto se han usado para almacenar los campos extraídos en una base de datos.
\end{description}

\imagen{scrapy_architecture}{Arquitectura que sigue Scrapy y sus diferentes componentes}

\subsection{\english{Spiders} utilizados}
El primer \english{spider} se encarga de extraer los enlaces a los productos. Para ello partimos de la página principal donde se encuentran todos los artículos a extraer. Cabe destacar que no existe un enlace que de acceso a una página web que contenga los productos orientados a ambos sexos, por lo que en caso de querer artículos para hombres y mujeres se necesita una url para cada sexo.

Este \english{spider} además genera un archivo JSON que contiene todos los enlaces que ha extraido. Este archivo será utilizado como punto de partida de los siguientes \english{spiders} para la extracción de los campos de cada artículo.

\imagen{principal-hombres2}{Productos a extraer}

El siguiente \english{spider} es el encargado de extraer la mayoría de campos de un producto. Es necesario destacar los diferentes métodos utilizados para la extracción de estos campos:

\begin{itemize}
    \item Expresiones regulares: El campo del ASIN se extrae directamente del url con ayuda de una expresión regular.
    \item La mayoría de campos se han extraido con la ayuda de expresiones XPath\footnote{\url{https://es.wikipedia.org/wiki/XPath}}, que consisten en parsear el código HTML de la web e identificar el campo que se desea extraer.
    \item Para las imágenes ha sido necesaria la utilización de una librería externa llamada js2xml\footnote{\url{https://github.com/scrapinghub/js2xml}}. Esta librería facilita la extracción de elementos embebidos en código JavaScript usando XPath.
\end{itemize}

\imagen{campos_1}{Posición de algunos de los campos a extraer}

Por último, el tercer \english{spider} del que hacemos uso es el encargado de extraer los comentarios de cada producto. Para ello usamos una expresión XPath sobre una página cuyo url creamos a partir del ASIN de cada artículo.

\imagen{comentarios2}{Comentarios sobre un producto}

En cuanto al \english{pipeline} de nuestro proyecto de Scrapy, cada vez que se extrae un \english{item}, éste es almacenado en una base de datos. El \english{pipeline} se encarga de conectar con la base de datos, crear las tablas necesarias con sus respectivos campos, y cerrar la conexión una vez se han almacenado todos los productos.

\imagen{proyecto_scrapy}{Proceso de \english{web scraping} del proyecto}

\section{Etiquetado manual de las imágenes}

Antes de etiquetar automáticamente las imágenes, es necesario entrenar el clasificador para que éste sea capaz de identificar cada imagen. Para ello se ha tenido que crear un conjunto de datos etiquetado manualmente que sirva de entrenamiento para el clasificador.

Ese conjunto de datos etiquetados o \english{dataset}, se ha generado con la ayuda de la web \english{Dataturks}\footnote{\url{https://dataturks.com}}.

El funcionamiento de esta herramienta es sencillo: basta con crear una lista con los enlaces a las imágenes que se desean etiquetar e indicar las posibles etiquetas que una imagen puede tener. A partir de aquí, la herramienta descarga y muestra secuencialmente todas las imágenes y se ha de decidir que etiqueta corresponde a cada una.

\imagen{dataturks}{Captura de pantalla durante el proceso de etiquetado manual}

Una vez se han etiquetado todas las imágenes, se puede exportar el conjunto de imágenes como un archivo JSON que asocia cada enlace con las etiquetas seleccionadas.

Para este proyecto se han creado dos \english{datasets}. En el primero se diferencian las imágenes en las que aparece un modelo y en las que no, el segundo \english{dataset} parte del anterior e identifica en cuales de las imágenes en las que si hay modelo, a éste se le ve la cara.

\subsection{Problemas encontrados}

El primero de los problemas sufridos relacionados con el etiquetado de imágenes fue precisamente decidir qué etiquetas seleccionar. En un principio la idea era identificar si aparecía un maniquí o no, o incluso si en la imagen no aparecía el producto anunciado. Estas ideas se desecharon ya que había demasiados casos en los que no era demasiado claro qué etiqueta correspondía a cada producto.

Otro problema encontrado ha sido a la hora de identificar si al modelo se le ve la cara o no. Hay muchos casos en los que la cara del modelo está cortada y solo es visible parcialmente. También hay casos en los que el modelo está de espaldas y mirando hacia un lado y solo se intuye la cara de éste. Lo primero que se pensó para estos casos fue un algoritmo que identificase las caras parcialmente, y aunque se han encontrado varios recursos que parecen dar respuesta a este problema, al final se optó por etiquetar las imágenes de forma que si parte de la cara es visible, se clasificaría como que si hay cara.

Algunos de los recursos que se han utilizado para una posible detección parcial de la cara han sido los siquientes:
\begin{itemize}
    \item \english{Partial Face Recognition}: \url{http://www.faceforensics.com/PartialFaceRecog.aspx}
    \item \english{Partial Face Detection in the Mobile Domain}: \url{https://arxiv.org/pdf/1704.02117.pdf}
    \item \english{Dynamic Feature Learning for Partial Face Recognition}: \url{http://openaccess.thecvf.com/content_cvpr_2018/papers/He_Dynamic_Feature_Learning_CVPR_2018_paper.pdf}
    \item \english{A robust face detector under partial occlusion - IEEE Conference Publication}: \url{https://ieeexplore.ieee.org/document/1418825}
    \item \english{Dealing with Inaccurate Face Detection for Automatic Gender Recognition with Partially Occluded Faces}: \url{http://repositori.uji.es/xmlui/bitstream/handle/10234/23941/34130.pdf}
\end{itemize}

%\imagen{cara_parcial}{Modelo con la cara parcialmente visible}
\imagen{cara_espalda_parcial}{Ejemplos de modelos con la cara parcialmente visible}

\section{Etiquetado automático de las imágenes}

Para el etiquetado automático de imágenes se ha decidido usar Keras. Existen otros \english{frameworks} que se han tenido en cuenta para este proyecto como pueden ser Pytorch\footnote{\url{https://pytorch.org/}} o Caffe\footnote{\url{https://caffe.berkeleyvision.org/}}, pero se ha usado Keras por varias razones:

\begin{description}
    \item[Experiencia con este \english{framework}:] Ya había trabajado con anterioridad con Keras para clasificación de imágenes.
    \item[Facilidad de uso:] Después de leer e informarme sobre otros \english{frameworks}, Keras me ha parecido el más facil de aprender y de poner en marcha.
    \item[Flexibilidad:] Keras se integra con TensorFlow\footnote{\url{https://www.tensorflow.org/guide/keras}}, es por esto que permite implementar cualquier cosa implemetada con TensorFlow como base.
    \item[Compatibilidad:] Keras es compatible con Google Cloud\footnote{\url{https://cloud.google.com/tpu/}} y Tarjetas gráficas NVIDIA\footnote{\url{https://developer.nvidia.com/deep-learning}}, gracias a ello funciona de forma nativa con Google Colab.
\end{description}

\subsection{Validación cruzada}
Para el entrenamiento de estos clasificadores se ha usado la técnica de validación cruzada. Esta técnica consiste en dividir el \english{dataset} en dos conjuntos complementarios y realizar el análisis sobre un subconjunto para luego validarlo con los datos del otro subconjunto. La principal ventaja de este método es que es muy rápido de computar.

Existen variaciones de este método, en el que se suceden varias iteraciones, obteniendo así una mayor precisión. En este caso se ha optado por el método de validación cruzada aleatoria ya que la precisión era aceptable y Keras dispone de una función que facilita este procedimiento \cite{wiki:cruzada}. 

\imagen{Random_cross_validation}{Validación cruzada aleatoria con $k$ iteraciones}
%\imagen{test_train}{Diagrama de conjuntos de entrenamiento y validación}

Para los clasificadores entrenados en este proyecto, se ha hecho una división del 75\% para el conjunto de entrenamiento y un 25\% para el conjunto de test. En el caso del clasificador que identifica si aparece un modelo o no, se obtiene una precisión de aproximadamente un 93\%. En cambio, para el clasificador que indica si parece la cara del modelo, la precisión es de únicamente un 80\%. La diferencia de precisión entre ambos clasificadores se debe a la complejidad del problema, ya que incluso se han tenido dudas a la hora elegir qué etiquetas usar de forma manual.

\imagen{precision_cara}{Precisión del clasificador Cara/No cara}

El número de iteraciones para el entrenamiento ha sido de 10 para ambos clasificadores. Se ha probado con distintos números de iteraciones pero la precisión obtenida no variaba o decrecía.

\imagen{summary}{Resumen de las capas del clasificador de imágenes}

\subsection{Problemas encontrados}

El principal inconveniente que se ha encontrado en el proceso de clasificación está directamente relacionado con el problema encontrado a la hora de etiquetar las imágenes manualmente. Para el clasificador que identifica si la cara de un modelo es visible, se pensó en utilizar una función de la librería de OpenCV que detecta caras en una imagen con la utilización de \english{Haar Cascades}\footnote{\url{https://docs.opencv.org/3.4.1/d7/d8b/tutorial_py_face_detection.html}}. Los resultados parecían prometedores al principio ya que actúa con una precisión muy alta. Pero el principal problema que se encontró es que no detectaba las caras cortadas. Esto hizo que se desechara la idea y se pasase a entrenar dos clasificadores.

\imagen{opencv_error}{Detección errónea de caras con la librería OpenCV}

\section{Visualización y almacenamiento de todos los productos junto a sus campos}

En este punto, ya tendríamos los clasificadores funcionando y los campos de los productos extraídos. Los pasos siguientes se encargan de satisfacer los objetivos de almacenamiento y visualización de los artículos extraídos:

\begin{enumerate}
    \item Cargar los clasificadores previamente entrenados
    \item Importar todos los productos que se desea clasificar. Esto se hace a partir del archivo JSON que ha generado Dataturks combinado con el archivo JSON creado por Scrapy.
    \item Clasificamos la imagen  principal de cada producto.
    \item Conectamos con la base de datos y creamos una nueva tabla que incluirá las predicciones de cada imagen. Es necesario recordar que la base de datos hasta este momento contiene 3 tablas: Productos, Imágenes y Comentarios. Esta nueva tabla asocia cada predicción con el ASIN del artículo al que pertenece.
\end{enumerate}

Con esto la base de datos quedaría completada. Cabe destacar que se ha elegido crear una nueva tabla en vez de añadir el campo <<predicción>> a una tabla ya existente ya que en un futuro se tendría pensado clasificar mas aspectos de cada artículo como podrían ser los comentarios u otras características de las imágenes.

\subsection{Almacenamiento de los datos}

La base de datos utilizada para almacenar los diferentes productos no es muy compleja. Una vez finalizada la extracción y el procesamiento de los datos, cuenta con 4 tablas con el ASIN de cada artículo como identificador.

\begin{itemize}
    \item La primera tabla, la más general, contiene los diferentes campos de cada producto y el ASIN como clave primaria. Cada entrada pertenece a un artículo diferente.
    \item La tabla de las imágenes almacena los enlaces a cada imagen junto al ASIN del producto al que pertenece. Puede haber varias imágenes para un único producto.
    \item La tabla de los comentarios almacena comentarios asociados a cada artículo. Puede haber varios comentarios de cada producto.
    \item La tabla de predicciones asocia cada artículo con la predicción que el clasificador ha generado para la imagen principal de dicho producto. Cada entrada es de un producto diferente.
\end{itemize}

Por último, es necesario añadir el proceso de creación de un documento Excel a partir de la base de datos ya completada. Para esto se ha utilizado una librería de Python llamada Pandas\footnote{\url{https://pandas.pydata.org/pandas-docs/stable/index.html}}, la cual extrae la información de la base de datos y la exporta a un documento Excel creando una página por cada tabla existente en la base de datos.

\imagen{bd_diagram}{Diagrama de la base de datos utilizada.}

\capitulo{6}{Trabajos relacionados}

Aunque existen multitud de proyectos que hacen uso de Scrapy para hacer \english{web scraping}, y muchos otros que utilizan Keras para clasificar imágenes; dado que este proyecto ha sido pensado con unos objetivos concretos, es difícil encontrar otros trabajos que hagan uso de estas dos herramientas o tengan unos objetivos parecidos.

Aún así son varios los artículos o proyectos que han servido de referencia:

\subsubsection{\english{How to use Scrapy to scrape Amazon.com and save to sqlite}}
\sloppy
\url{https://sunshineafternoonweb.wordpress.com/2017/01/30/how-to-use-scrapy-to-scrape-amazon-com-and-save-to-sqlite/}

\subsubsection{amazon-scrapy}
\url{https://github.com/dynamohuang/amazon-scrapy}

\subsubsection{Image Recognition with Keras: Convolutional Neural Networks}
\sloppy
\href{https://www.michael-grogan.com/image-recognition-with-keras-convolutional-neural-networks/}{https://www.michael-grogan.com/}
\linebreak
\url{image-recognition-with-keras-convolutional-neural-networks/}

\subsubsection{text-detector-OCR}
Mi propio trabajo realizado para la asignatura Minería de datos ha servido como referencia para multitud de puntos en el proceso de clasificación de imágenes usando Keras.\linebreak
\url{https://github.com/daniarnaizg/text-detector-ocr}

\capitulo{7}{Conclusiones y Líneas de trabajo futuras}

\section{Conclusiones generales}
En primer lugar y en relación a los principales requisitos y objetivos del proyecto, creo que se han completado de forma satisfactoria ofreciendo un forma de extraer información de los productos deseados en \english{Amazon.com} para más tarde clasificar las imágenes principales de cada artículo y almacenar toda esta información tanto en una base de datos como en un documento excel para su posterior uso en la realización de distintos estudios y cálculos.

En cuanto a las dificultades encontradas en la realización del proyecto, encontramos algunas en la parte de \english{web scraping} como pueden ser el cambio del código html de la página que ha hecho cambiar la forma de extraer los diferentes campos en diversas ocasiones, o el bloqueo temporal que la propia web ponía a mi dirección IP a la hora de extraer grandes cantidades de productos. En la parte del clasificador de imágenes también han surgido algunos inconvenientes, el principal problema encontrado está relacionado con la detección parcial de las caras de los modelos, lo cual ha significado probar diversos métodos y técnicas hasta encontrar la que mejor resultados ofrecía.

De forma más general, otro inconveniente encontrado ha sido la forma de unir las diferentes partes que conforman el proyecto. Aunque casi la totalidad del código es Python, algunas partes se han realizado en \english{notebooks} de Google Colab y otras dentro de la estructura que ofrece Scrapy. Es por esto que se ha decidido plantear el proyecto de forma modular creando una diferenciación entre cada una de las partes que lo conforman: \english{Web scraping}, clasificación de imágenes y Almacenamiento de datos.

A nivel personal, creo que ha sido una experiencia enriquecedora y positiva en varios sentidos, el primero puramente académico ya que he necesitado aprender y estudiar nuevos conceptos y técnicas además de reforzar lo ya tratado en la propia universidad. Y el segundo, ya de cara al futuro, con la puesta en práctica de metodologías y técnicas de investigación y de toma de decisiones.

\section{Líneas de trabajo futuras}

A continuación se listan algunas posibles mejoras o ampliaciones que podrían incorporarse al proyecto en versiones futuras:

\begin{itemize}
    \item Integración de las distintas partes del proyecto en una interfaz de usuario para hacer mas sencilla la utilización de la herramienta.
    \item Clasificar y almacenar todas las imágenes asociadas a un artículo en vez de únicamente la principal.
    \item Hacer una clasificación de los comentarios o de las descripciones de los artículos. En esta versión únicamente se almacenan y se asocian con el producto al que pertenecen.
    \item Ampliar la clasificación de imágenes indicando mas características. Algunos ejemplos podrían ser la detección de algunas partes del cuerpo, o qué colores aparecen en la imagen.
    \item Dotar a la herramienta de una forma de visualizar los datos finales de forma interactiva.
\end{itemize}


\bibliographystyle{plain}
\bibliography{bibliografia}

\end{document}
